\subsection{বিটের দৈর্ঘ্য, পরিসর ও ফল ব্যাখ্যা}
চিহ্নবিহীন এবং চিহ্নযুক্ত সংখ্যা একই বিটের pattern ব্যবহার করলেও তাদের মান ভিন্ন হতে পারে। $n$ বিটে চিহ্নবিহীন মান $0$ থেকে $2^n-1$। সাইন-ম্যাগনিটিউড ও ১-এর পরিপূরকে পরিসর $-(2^{n-1}-1)$ থেকে $+(2^{n-1}-1)$; দুটিতে শূন্যের দুটি প্রকাশ আছে। ২-এর পরিপূরকে পরিসর $-2^{n-1}$ থেকে $2^{n-1}-1$ এবং শূন্য একটিই।

২-এর পরিপূরকে সর্ববামের বিটের ওজন ঋণাত্মক:
\[V=-d_{n-1}2^{n-1}+\sum_{i=0}^{n-2}d_i2^i.\]
তাই সর্ববামের বিটকে সরিয়ে বাকি বিটের সাধারণ মান পড়লেই ঋণাত্মক সংখ্যার মান পাওয়া যায় না।

\chThreeExample{১}{একই pattern-এর চারটি অর্থ}
আট বিটের $11110110$ pattern-এর মান নির্ণয় করো।
\chThreeAnswer চিহ্নবিহীন হিসেবে ২৪৬; সাইন-ম্যাগনিটিউডে $-118$; ১-এর পরিপূরকে উল্টালে $00001001$, তাই $-9$; ২-এর পরিপূরকে $-128+64+32+16+4+2=-10$। প্রশ্নে প্রকাশপদ্ধতি উল্লেখ থাকা প্রয়োজন।

\chThreeExample{২}{ধনাত্মক থেকে ঋণাত্মক, তারপর যোগ}
আট বিটে $18+(-27)$ নির্ণয় করো।
\chThreeAnswer $18=00010010$, $27=00011011$। $-27=11100101$।
\[00010010+11100101=11110111.\]
ফল ঋণাত্মক; এর ২-এর পরিপূরক $00001001=9$। অতএব ফল $-9$।

\chThreeExample{৩}{দুটি ঋণাত্মক সংখ্যার যোগে বেশি বিট}
$(-86)+(-57)$-এর সঠিক ফল পেতে ১৬ বিট ব্যবহার করো।
\chThreeAnswer $86=0000000001010110$, তাই $-86=1111111110101010$। $57=0000000000111001$, তাই $-57=1111111111000111$।
\[\begin{array}{r}1111111110101010\\+1111111111000111\\\hline1\;1111111101110001\end{array}\]
অতিরিক্ত ক্যারি বাদ দিলে $1111111101110001$। এর পূরক নিয়ে মান $143$ পাওয়া যায়, তাই ফল $-143$। আট বিটের signed পরিসরে $-143$ নেই; এজন্য বেশি বিট নেওয়া হয়েছে।

\chThreeExample{৪}{ঋণাত্মক সংখ্যা থেকে ঋণাত্মক সংখ্যা বিয়োগ}
$(-14)-(-23)$।
\chThreeAnswer $-14=11110010$ এবং $+23=00010111$। কারণ $A-(-B)=A+B$।
\[11110010+00010111=1\;00001001\Rightarrow +9.\]

\chThreeExample{৫}{সাইন এক্সটেনশন}
চার বিটে $-3=1101$। আট বিটে একই মান রাখতে বামে সাইন বিট ১ পুনরাবৃত্ত করতে হয়: $11111101$। $00001101$ করলে মান $+13$ হয়ে যায়। ধনাত্মক সংখ্যায় বামে ০ বসে।

\chThreeExample{৬}{ক্যারি এবং signed overflow পৃথক}
চার বিটে $7+3$: $0111+0011=1010$। বাইরে ক্যারি নেই, কিন্তু signed ফল $-6$; প্রকৃত যোগফল ১০ চার বিটের signed পরিসরে নেই। তাই overflow হয়েছে। আবার $(-1)+1$: $1111+0001=1\;0000$; ক্যারি আছে, কিন্তু ফল ০ সঠিক এবং overflow নেই। একই চিহ্নের দুটি সংখ্যা যোগ করে বিপরীত চিহ্নের ফল পেলে signed overflow শনাক্ত হয়।

\smallhead{২-এর পরিপূরক কেন সার্কিট সহজ করে?}
বিয়োগ $A-B$-কে $A+\overline B+1$ হিসেবে করা যায়। একই অ্যাডারে $B$-এর বিটগুলো প্রয়োজনমতো উল্টে এবং প্রথম ক্যারি ইনপুটে ১ দিয়ে বিয়োগ সম্ভব। আলাদা যোগ ও বিয়োগের পুরো সার্কিট দরকার হয় না। এই যুক্তিই পরিপূরক পদ্ধতির প্রধান শিক্ষণীয় সুবিধা।

% Mahbub teaching-type follow-up, 2026-10-14
\chThreeExample{৭}{ঋণাত্মক সংখ্যার দুই ধরনের বিয়োগ}
আট বিটের ২-এর পরিপূরকে $(-31)-(+18)$ এবং $(-31)-(-18)$ নির্ণয় করো।
\chThreeAnswer $-31=11100001$, $+18=00010010$, $-18=11101110$।
\[\begin{aligned}
(-31)-(+18)&:\quad11100001+11101110=1\;11001111,\\
(-31)-(-18)&:\quad11100001+00010010=11110011.
\end{aligned}\]
প্রথম ফলে অতিরিক্ত carry বাদ দিয়ে 11001111-এর পরিপূরকের মান ৪৯; ফল $-49$। দ্বিতীয়টিতে 11110011-এর পরিপূরকের মান ১৩; ফল $-13$। দুটি ঋণাত্মক সংখ্যার বিয়োগের ফল সবসময় ধনাত্মক হয় না; তাদের মানের তুলনা প্রয়োজন।

\chThreeExample{৮}{সমান সংখ্যার বিয়োগে শূন্য}
আট বিটে পরিপূরক যোগের সাহায্যে $46-46$ দেখাও।
\chThreeAnswer $46=00101110$; এর ২-এর পরিপূরক $11010010$।
\[00101110+11010010=1\;00000000.\]
বাইরের carry বাদ দিলে ফল ০। শূন্য ফলকে ঋণাত্মক বলার প্রয়োজন নেই; ২-এর পরিপূরকে শূন্যের একটিমাত্র প্রকাশ।
